\section{The Testbed and Its Discipline}
\label{sec:testbed}

The findings come from \texttt{@cluesurf/vibe}, a finite, seeded, reproducible simulator
\citep{pollard2026vibetest}. Its purpose is to keep the theory honest by turning each claim into a measurement that
could come out wrong.

\paragraph{Finite and seeded.}
Every run is finite and deterministic given a seed. There is no hidden state and no wall-clock
randomness. Re-running a finding reproduces it exactly. This matters because the model's claims are
quantitative, and a quantitative claim is only as trustworthy as its reproducibility.

\paragraph{The discreteness principle.}
Real numbers appear only as measured outputs, never as the substrate. The base carries
three-valued states and integer counts. Continuous quantities, coordinates, eigenvalues,
dimensions, temperatures, are computed from the discrete data, in the spirit of a cellular
automaton whose tally of neighbours is an integer. When a continuous quantity such as a Poincar\'e
coordinate is used, it is a derived embedding of a discrete graph, not a primitive.

\paragraph{Known-answer tests.}
Each finding is paired with a test that checks it against an independently known value: a
classification from the mathematics literature, a textbook physics limit, an analytic result, or a
self-consistency identity. The suite currently holds one hundred three such checks, all passing.
Several early versions of findings failed these tests in plausible-looking ways and were corrected.
The tests are the reason the numbers in this paper can be trusted.

\paragraph{How a finding is reported.}
Each numbered problem, written P1, P2, and so on, is one experiment. Its code lives in
\texttt{code/experiment/}, its result note in \texttt{note/experiment/results/}, and its status in
\texttt{note/questions/}. In this paper we group the findings by theme rather than by number, and
for each we give the method in one or two sentences, the result, and an honest status: demonstrated,
a first rung, or open. The full per-finding scoreboard is Section~\ref{sec:scoreboard}.

\paragraph{The committed model.}
One configuration is the default everywhere, the committed model: a growing hyperbolic causal mesh,
ternary tones $\{-1,0,+1\}$, ternary fills on the notes, a signed-majority update, asynchronous
local scheduling, and net-positive growth. A small constructor writes it, and the capstone of
Section~\ref{sec:capstone} runs it end to end, reading the key emergent structures off one
instantiation.
